% !TeX root = ../main.tex
% ---------------------------------------------------------------------------
% In-situ Stress: Estimation and Measurement
% Scope extension: deepens the in-situ stress treatment in Chapter 2 ( Chapter~\ref{ch:stress-intact-rock}).
% Key sources: Amadei & Stephansson; Zoback, World Stress Map; Sheorey (1994);
% ISRM Suggested Methods for stress determination.
% ---------------------------------------------------------------------------
\chapter{In-situ Stress: Estimation and Measurement}
\label{ch:in-situ-stress}


\section{The in-situ stress field}
\label{sec:the-in-situ-stress-field}

\subsection{Why in-situ stress matters}
\label{sec:why-stress-matters}
\scopenote{Stress sets the demand side of every underground design check. Its
influence on failure mode, excavation shape, orientation and support demand.}

\subsection{Components of the in-situ stress field}
\label{sec:stress-components}
\scopenote{Gravitational loading, plate-boundary tectonic stress, topographic
effects in valleys and mountains, and residual stress locked in during
geological history.}

\subsection{Regional stress fields and the Australian setting}
\label{sec:regional-stress}
\scopenote{The World Stress Map, the Australian Stress Map, and the unusual
character of the Australian intraplate stress field --- high horizontal
stresses with regionally variable orientation.}

\section{Vertical stress and the $k$-ratio}
\label{sec:k-ratio}
\scopenote{$\sv = \gamma z$ and its exceptions; the Hoek--Brown $k$ envelope
with depth; Sheorey's thermo-elastic model as a physically based estimate.}

\section{Measurement methods}
\label{sec:measurement-methods}

\subsection{Overcoring methods}
\label{sec:overcoring}
\scopenote{CSIRO HI cell, USBM gauge and doorstopper. Full stress tensor
recovery, required rock quality, depth limits and typical scatter.}

\subsection{Hydraulic methods}
\label{sec:hydraulic-fracturing}
\scopenote{Hydraulic fracturing and HTPF. What is measured directly, what is
inferred, and the assumption that a principal stress is vertical.}

\subsection{Flat jack, borehole slotter and relief methods}
\label{sec:flat-jack}
\scopenote{Near-surface and in-excavation techniques, and their sensitivity to
the disturbed zone they sit within.}

\section{Indirect evidence and back-analysis}
\label{sec:indirect-evidence-and-back-analysis}

\subsection{Stress indicators from boreholes and geology}
\label{sec:stress-indicators}
\scopenote{Borehole breakout and drilling-induced fractures, core discing,
focal mechanisms and fault slip data as orientation evidence when direct
measurement is unavailable.}

\subsection{Back-analysis of stress from observed behaviour}
\label{sec:stress-back-analysis}
\scopenote{Inferring the stress field from convergence, overbreak pattern and
observed failure geometry during construction.}

\section{Designing under stress uncertainty}
\label{sec:designing-under-stress-uncertainty}

\subsection{Variability, and designing to a stress envelope}
\label{sec:stress-envelope}
\scopenote{Measurement scatter of 20--50\,\% is normal. Constructing upper and
lower bound cases and testing the design against both, rather than adopting a
single best estimate.}

\subsection{Worked example: stress estimation for a deep tunnel}
\label{sec:stress-worked-example}
\scopenote{Estimate $\sv$, $\sh$ and orientation for a deep alignment from
depth, topography and regional data; compare with the Snowy 2.0 case in
Chapter~\ref{ch:ccm}.}
